% === Begin Label Data ===
% ID: 583
% 难度星级: 
% 标签: 
% 备注: 
% 组卷引用次数: 0
% === End  Label Data ===

\begin{problem}{2023}{M}{2023年高考模拟试题(Fiddie卷)}{18}{解三角形}
十字测天仪广泛应用于欧洲中世纪晚期的航海领域，主要用于测量太阳等星体的方位，便于船员确定位置。如图1所示，十字测天仪由杆 $ AB $ 和横档 $ CD $ 构成，并且 $ E $ 是 $ CD $ 的中点，横档与杆垂直且可在杆上滑动。十字测天仪的使用方法如下：如图2，手持十字测天仪，使得眼睛可以从 $ A $ 点观察。滑动横档 $ CD $ 使得 $ A $、$ C $ 在同一水平面上，并且眼睛恰好能观察到太阳，此时视线恰好经过点 $ D $，$ DE $ 的影子恰好是 $ AE $.然后，通过测量 $ AE $ 的长度，可计算出视线和水平面的夹角 $ \angle CAD $（称为太阳高度角），最后通过查阅地图来确定船员所在的位置。

（1）在某次测量中，$ AE = 40 $，横档的长度为 $ 20 $，求太阳高度角的正弦值。

（2）在杆 $ AB $ 上有两点 $ A_1 $、$ A_2 $ 满足 $ AA_1 = \dfrac{1}{2}AA_2 $.当横档 $ CD $ 的中点 $ E $ 位于 $ A_i $ 时，记太阳高度角为 $ \alpha_i $（$ i = 1,2 $），其中 $ \alpha_1 $、$ \alpha_2 $ 都是锐角。证明：$ \alpha_1 < 2\alpha_2 $.
\end{problem}

\begin{answer}

\end{answer}

\begin{solutions}

\end{solutions}